\section{The genuine Gauss point on each literal original test configuration}

\begin{definition}\label{def:native-test-configuration-gauss-point}
\lean{EnvelopeProperty.NativeTestConfiguration.TestConfiguration.unitGaussPoint, EnvelopeProperty.NativeTestConfiguration.TestConfiguration.totalGaussPoint}
\leanok
Let $C$ be a literal original test configuration for $X$ over a field $k$,
let $v$ be a valuation of $k$, let $x$ be an original native point of $X$
extending $v$, and let $r>0$. Construct the actual completed-field Laurent
Gauss product point of $X\times\mathbf G_m$, map it through the inverse
original noncentral-open trivialization of $C$, and then through the
original open inclusion into the original total scheme. Denote the
resulting actual native total-space point by $y_{C,x,r}$.
\uses{thm:native-point-laurent-gauss-projections,thm:native-test-configuration-unit-fibre}
\end{definition}
\begin{proof}
The preceding construction supplies the actual whole product point and
its original coefficient and parameter projections from $H(x)$.
The literal original source action and fibre-one isomorphism supply the
actual whole scheme isomorphism between the original noncentral open
and $X\times\mathbf G_m$. Its inverse and the actual open inclusion give
the asserted actual total-space point. No selected lift or projection
matching data are added as hypotheses.
\end{proof}

\begin{theorem}\label{thm:native-test-configuration-gauss-point-base}
\lean{EnvelopeProperty.NativeTestConfiguration.TestConfiguration.unitGaussPoint_originalProjection, EnvelopeProperty.NativeTestConfiguration.TestConfiguration.totalGaussPoint_projection, EnvelopeProperty.NativeTestConfiguration.TestConfiguration.relativeGaussPoint}
\leanok
The whole polynomial-base image of $y_{C,x,r}$ equals the original
polynomial Gauss point of $v$ at radius $r$. Thus $y_{C,x,r}$ is an actual
relative valued point of the original projection $C\to\mathbf A^1_k$
with precisely that prescribed valuation. On the original noncentral
open its whole $X$ projection recovers exactly $x$.
\uses{def:native-test-configuration-gauss-point,thm:native-point-laurent-gauss-projections}
\end{theorem}
\begin{proof}
The whole original noncentral-open trivialization retains the actual
parameter morphism to the original affine line. Cancel its isomorphism
to identify the polynomial-base composite of the constructed point
with the original product parameter projection. The actual completed-field
Gauss point projects to the entire original polynomial Gauss point by the
preceding theorem. This proves the prescribed-base equation, hence the
relative valued-point assertion. Cancelling the same isomorphism in the
original first projection gives the entire original $X$-point equality.
These arguments compare actual native valued points, not only their
supports, valuations of the parameter, or chosen affine coordinates.
\end{proof}
